% ---
% id: meat-core/cor:composed-security
% kind: cor
% title: "Security of the protocol and its net key"
% summary: >-
%   Secrecy and correctness compose to (ε_PA + ε_EV)-security with honest
%   authentication; over an authentication scheme within ε_auth that consumes
%   k_auth bits of pre-shared key, the protocol is (ε_PA + ε_EV +
%   2ε_auth)-secure; when both parties output a key, its net key gain is
%   ℓ − k_auth bits.
% user_statement: >-
%   A run either aborts or outputs a key of ℓ bits that is a perfect key up
%   to the security parameter ε: identical for both parties, uniform and
%   independent of all the attacker holds, public messages included. One
%   party may abort while the other gets a key. No test, whatever the
%   attacker knows, tells the run's output from an ideal protocol's with
%   advantage above ε = ε_PA + ε_EV + 2ε_auth. It holds against every attack,
%   of any memory, if every premise in the run record holds, [[A-AUTH]],
%   [[A-CONFORM]] and [[A-FIXED]] included. A first-order run also needs its
%   bound to hold, computed h_lb ≤ h_low, which no premise covers
%   ([[meat-core/frank-wolfe/rem:fo-arithmetic]]). ε averages over aborted
%   runs too. The criterion is composable. Separate executions provably add ε
%   ([[cor:separate-executions]]). Not every later use of the key is covered
%   ([[rem:repeated-executions]]). A run where both get a key nets ℓ − k_auth
%   bits; an abort uses up to k_auth. Abort rates:
%   [[rem:completeness-scope]].
% prominence: main
% status: implemented
% uses: [thm:keylength, lem:ev-correctness, lem:memory-lift, lem:honest-attacks, def:postprocessing-model, def:authenticated-channel]
% mentions: [def:execution-end]
% assumes: [A-AUTH, A-FIXED, A-CONFORM]
% refs:
%   - "tupkary2025survey, Sec. 2.3, Def. 2.1, Eq. (1)"
%   - "tupkary2025survey, Sec. 2.3.1, Thm. 2.1"
%   - "portmann2022security, Sec. III.B.3, Thm. 2"
%   - "portmann2022security, Sec. III.B.1"
%   - "portmann2022security, Sec. II.E, Def. 1"
%   - "portmann2022security, Sec. II.F, Eqs. (4)-(6)"
%   - "portmann2022security, Sec. VII.D"
%   - "tupkary2026decoy, Sec. 3.3, Def. 3.8, Def. 3.9, Eqs. (13)-(17), Eq. (23)"
%   - "tupkary2026decoy, Sec. 5, Thm. 5.1 and the paragraph following it"
%   - "tupkary2026decoy, Sec. 5, Eqs. (37)-(38)"
%   - "tupkary2026decoy, Sec. 4.1, Protocol 1, steps 2-6"
%   - "tupkary2026decoy, Sec. 4, footnote 8"
%   - "tupkary2025survey, Sec. 5.6.1"
% provenance: newly-derived
% implements: certiqs_key.meat_core.finitekey.keylength:composed_security
% lean:
% ---

Secrecy (Theorem~\ref{thm:keylength}) and correctness
(Lemma~\ref{lem:ev-correctness}) are the two halves of the security of the
fixed-length protocol in the setting where authentication behaves honestly.
A real authentication scheme can abort asymmetrically and has its own error;
the authentication post-processing reduces the first to the honest setting
\cite[Sec.~5, Thm.~5.1]{tupkary2026decoy}, and the scheme's error enters the
security parameter, here through the triangle inequality, which charges it
twice. That it enters is the motivation of
\cite[Sec.~4, footnote~8]{tupkary2026decoy} and
\cite[Sec.~5.6.1]{tupkary2025survey}; (iii) is derived without them.

\begin{corollary}\label{cor:composed-security}
Let the hypotheses of Theorem~\ref{thm:keylength}, the run as in
Lemma~\ref{lem:memory-lift}\,(a) and (H-CHAIN) among them, hold at every
attack in the sense of Lemma~\ref{lem:memory-lift}\,(b) that is all-at-once
and has finite registers, with $\ell\ge1$ satisfying \eqref{eq:keylength},
A-SEQ with its clause (v) among them; let those of
Lemma~\ref{lem:ev-correctness} hold; and put
$\tau_{K_AK_B}\coloneqq2^{-\ell}\sum_{k\in\{0,1\}^{\ell}}
\dyad{k}_{K_A}\otimes\dyad{k}_{K_B}$.
\begin{enumerate}[label=(\roman*),leftmargin=2.4em]
\item In the \tref{honest-authentication setting}, at every output state the
  setting admits, whatever the dimension of the adversary's registers and the
  order in which the source releases its signals, the protocol is
  $\epa$-secret on the accepted state as Theorem~\ref{thm:keylength} states
  it, and $(\epa+\eev)$-secure in the sense of
  \cite[Sec.~2.3, Def.~2.1, Eq.~(1)]{tupkary2025survey}:
  \begin{equation}\label{eq:composed-honest}
    \frac12\Big\|\rho_{K_AK_B\,E_{\mathrm{tot}}\wedge\Omega_{\mathrm{acc}}}
    -\tau_{K_AK_B}\otimes
    \rho_{E_{\mathrm{tot}}\wedge\Omega_{\mathrm{acc}}}\Big\|_1
    \le\epa+\eev ,
  \end{equation}
  and on the complement of $\Omega_{\mathrm{acc}}$ both parties output
  $\perp$.
\item Run with the \tref{authentication post-processing}
  (Definition~\ref{def:authenticated-channel}\,(c)) over the
  \tref{authenticated classical channel}
  (Definition~\ref{def:authenticated-channel}\,(a)), the protocol is
  $(\epa+\eev)$-secure in the sense of
  \cite[Sec.~3.3, Def.~3.8]{tupkary2026decoy}, which admits asymmetric
  aborts, for every output state that channel admits.
\item Let Assumptions~\ref{ass:auth}, \ref{ass:fixed}
  and~\ref{ass:conform} hold, with $\varepsilon_{\mathrm{auth}}$ and
  $k_{\mathrm{auth}}$ the parameters of Assumption~\ref{ass:auth}. Run with
  the authentication post-processing over the classical channel of
  Assumption~\ref{ass:auth}, the protocol is
  $(\epa+\eev+2\varepsilon_{\mathrm{auth}})$-secure in the sense of
  \cite[Sec.~3.3, Def.~3.8]{tupkary2026decoy}: for every attack on that run,
  with output state $\rho'$ on the key registers and the adversary's final
  information,
  \begin{equation}\label{eq:composed-auth}
    \frac12\big\|\rho'-\mathcal R_{\mathrm{ideal}}(\rho')\big\|_1
    \le\epa+\eev+2\varepsilon_{\mathrm{auth}},
  \end{equation}
  with $\mathcal R_{\mathrm{ideal}}$ the map of
  \cite[Sec.~3.3, Eqs.~(13)--(17)]{tupkary2026decoy}. It consumes at most
  $k_{\mathrm{auth}}$ bits of the pre-shared key, and when both parties
  output an $\ell$-bit key its \emph{net key gain} is
  $\ell_{\mathrm{net}}\coloneqq\ell-k_{\mathrm{auth}}$ bits; for
  $\ell_{\mathrm{net}}\le0$ the run yields no net key.
\end{enumerate}
\end{corollary}

\begin{proof}
\emph{(i).} We prove the two bounds at an all-at-once attack with finite
registers first, and then at every output state of the setting.

\emph{At an all-at-once attack with finite registers.} Theorem~\ref{thm:keylength}
gives the secrecy there. The rest is the statement that correctness and secrecy together imply
security \cite[Sec.~2.3.1, Thm.~2.1]{tupkary2025survey}, which is
\cite[Sec.~III.B.3, Thm.~2]{portmann2022security}; we give the argument in
this contract's notation. In the honest-authentication setting both parties
learn $\Omega_{\mathrm{acc}}$, a function of the announced $\Chat_1^n$ and
$F_{\mathrm{acc}}$, so on its complement both output $\perp$, where the real
and the ideal state coincide; it remains to bound
\eqref{eq:composed-honest}. Write
$\sigma\coloneqq\rho_{K\hat K E_{\mathrm{tot}}\wedge\Omega_{\mathrm{acc}}}
=\sum_{k,k'}\dyad{k}_K\otimes\dyad{k'}_{\hat K}\otimes\sigma^{k,k'}$, with
$\sigma^{k,k'}\ge0$ on $E_{\mathrm{tot}}$, and let $C$ be the channel on $K$
that measures $K$ in its classical basis and writes the outcome into $K$ and
into a copy $\hat K$, $C(X)\coloneqq\sum_k\bra{k}X\ket{k}\,
\dyad{k}_K\otimes\dyad{k}_{\hat K}$. Put
$\sigma'\coloneqq(C\otimes\idmap_{E_{\mathrm{tot}}})
(\rho_{K E_{\mathrm{tot}}\wedge\Omega_{\mathrm{acc}}})
=\sum_{k}\dyad{k}_K\otimes\dyad{k}_{\hat K}\otimes\sum_{k'}\sigma^{k,k'}$.
The difference $\sigma-\sigma'$ is block diagonal in $k$, the terms with
$k'=k$ cancel, and in block $k$ it is
$\sum_{k'\ne k}\dyad{k'}_{\hat K}\otimes\sigma^{k,k'}
-\dyad{k}_{\hat K}\otimes\sum_{k'\ne k}\sigma^{k,k'}$, whose trace norm is
$2\sum_{k'\ne k}\tr\sigma^{k,k'}$ because every $\sigma^{k,k'}$ is positive.
Hence
$\frac12\|\sigma-\sigma'\|_1=\Pr[\Omega_{\mathrm{acc}}\wedge K\ne\hat K]
\le\eev$ by Lemma~\ref{lem:ev-correctness}. Since $C$ maps
$\idon{K}/2^{\ell}$ to $\tau_{K\hat K}$ and a channel does not increase the
trace norm,
$\frac12\|\sigma'-\tau_{K\hat K}\otimes
\rho_{E_{\mathrm{tot}}\wedge\Omega_{\mathrm{acc}}}\|_1\le\epa$ by
Theorem~\ref{thm:keylength}. The triangle inequality gives
\eqref{eq:composed-honest}, read for $(K_A,K_B)$, which is $(K,\hat K)$ in
the order the key holder fixes; $\tau$ is symmetric under exchanging the two
registers.

\emph{At every output state of the setting.} With $\Pi$ the projection onto
$\Omega_{\mathrm{acc}}$, a function of the public messages, the secrecy
distance of Theorem~\ref{thm:keylength} is the criterion $f$ of
Lemma~\ref{lem:memory-lift}\,(ii) for $L_1(\sigma)$ the selection
$\Pi\sigma\Pi$ followed by the partial trace over every register other than
$K$ and the public registers of $E_{\mathrm{tot}}$, and
$L_2(\sigma)\coloneqq\frac{\idon{K}}{2^{\ell}}\otimes\tr_KL_1(\sigma)$; the
left side of \eqref{eq:composed-honest} is the same with $K_AK_B$ in place of
$K$ and $\tau_{K_AK_B}$ in place of $\idon{K}/2^{\ell}$. Each $L_t$ is
completely positive, does not increase the trace, keeps the public messages
and acts on the other registers given their value, as part (iii) of that
lemma asks, being composed of the selection by $\Pi$, which is diagonal in
the public messages, partial traces over other registers and the appending
of a state. By the first part both bounds hold at every all-at-once attack
with finite registers, so by Lemma~\ref{lem:memory-lift}\,(ii) at every attack
in the sense of its (b). By Lemma~\ref{lem:honest-attacks}, whose premises
A-SEQ, its clause (v) included, A-LAB at both parties, A-CONFORM and
A-ALICE-MEAS\,(i) where no \tref{source-replacement certificate} is demanded
are among those of Theorem~\ref{thm:keylength}, every output state of the
honest-authentication setting is $(\idmap\otimes\Phi)(\rho)$ for the output
state $\rho$ of such an attack and a channel $\Phi$ of the adversary given
the public messages, and by Lemma~\ref{lem:memory-lift}\,(iii) neither
criterion increases under $\Phi$. That both parties output $\perp$ off
$\Omega_{\mathrm{acc}}$ holds at every output state, being the run's
processing of its public messages.

\emph{(ii).} Write $\widetilde{\mathcal P}$ for the protocol of the
\tref{fixed-length post-processing model} and
$\mathcal P\coloneqq\mathcal P_{\mathrm{APP}}\circ\widetilde{\mathcal P}$ for
it followed by the authentication post-processing. The model is an instance
of the protocol of \cite[Sec.~4.1, Protocol~1, steps~2--5]{tupkary2026decoy}
with the key holder as Alice: the length and error-correction decisions are
$\ell(\hat c_1^n)=\ell\cdot\mathbf 1[\operatorname{freq}_{\hat c_1^n}\in
\acSet]$ and the constant $\lEC$, so the messages that carry them are public
functions of $\Chat_1^n$; error correction is one-way with at most
$2^{\lEC}$ transcripts; error verification and privacy amplification are
those steps with the hash lengths of
Definition~\ref{def:postprocessing-model}. By (i),
$\widetilde{\mathcal P}$ is $(\epa+\eev)$-secure, for every output state of
the honest-authentication setting, in the sense of
\cite[Sec.~3.3, Def.~3.9]{tupkary2026decoy} at fixed length
\cite[Eq.~(23)]{tupkary2026decoy}: its key lengths always coincide, and
$E_{\mathrm{tot}}$ holds every public message of $\widetilde{\mathcal P}$
(the announcements, the acceptance values, the seeds, $L$ and $G$) together
with the adversary's quantum memory. That is the hypothesis of
\cite[Sec.~5, Thm.~5.1]{tupkary2026decoy}, whose honest-authentication setting
is this contract's and whose practical authentication setting is the
authenticated classical channel of
Definition~\ref{def:authenticated-channel}\,(a), properties (A1) and (A2)
being its Eqs.~(37)--(38). The theorem concludes $(\epa+\eev)$-security of
$\mathcal P$ in the sense of Def.~3.8 for every output state of that channel;
by the paragraph following it, it applies to an arbitrary protocol followed
by the authentication post-processing of that work, which is
Definition~\ref{def:authenticated-channel}\,(c): step~6 of
\cite[Sec.~4.1, Protocol~1]{tupkary2026decoy} together with the rule of its
preamble that a party receiving an authentication failure aborts, with the
key holder in the role of Alice and the partner in that of Bob.

\emph{(iii).} By Assumption~\ref{ass:conform} the run is $\mathcal P$: it
applies the acceptance test and the verification abort, sends no message
beyond those of the model and of the authentication post-processing,
outputs the $\ell$-bit keys the model computes or $\perp$, uses no key
before the \tref{end of an execution} (Definition~\ref{def:execution-end}),
and aborts on an authentication failure. By
Assumption~\ref{ass:fixed}, $\varepsilon_{\mathrm{auth}}$ and
$k_{\mathrm{auth}}$, with every other parameter, are fixed before the run,
and by Assumption~\ref{ass:auth}\,(i) the run uses one scheme, so the scheme
and the protocol are fixed systems. If the classical channel
is itself the authenticated classical channel
(Assumption~\ref{ass:auth}\,(i), second case), the run is one of those of
(ii), $\varepsilon_{\mathrm{auth}}=0$, and \eqref{eq:composed-auth} is (ii).
Otherwise fix an attack on the run over the scheme, with output state
$\rho'$. By Assumption~\ref{ass:auth}\,(iii) and
Definition~\ref{def:authenticated-channel}\,(b) there is a simulator $\sigma$
at the adversary's interface such that no distinguisher tells the scheme,
run on two identical copies of a uniformly distributed key independent of
everything else and the unauthenticated channel, from the authenticated classical channel
with $\sigma$ with advantage above $\varepsilon_0$
\cite[Sec.~II.E, Def.~1; Sec.~II.F, Eqs.~(4)--(6)]{portmann2022security}.
The run, the attack and any distinguisher after them form a channel applied
to the joint state, before the run, of the two parties' copies of the
pre-shared key with every other register and the adversary's information, together with fresh systems
independent of it; by Assumption~\ref{ass:auth}\,(iv) and because a channel
does not increase the trace distance, every final state against the scheme
run on the actual copies lies within $\delta$ of the one against the scheme
run on the ideal copies. So the scheme run on the actual copies is told from
the
authenticated classical channel with $\sigma$ with advantage at most
$\varepsilon_0+\delta\le\varepsilon_{\mathrm{auth}}$
(Assumption~\ref{ass:auth}\,(v)).
Consider the distinguisher that plays both parties of $\mathcal P$, every
quantum round included, and the given attack, which reaches the classical
channel only through the adversary's interface of the system the
distinguisher is given, and which then holds the parties' key registers and
the adversary's final information. Its final state is $\rho'$ against the
scheme and, against the authenticated channel with $\sigma$, the output
state $\rho$ of $\mathcal P$ under the attack that the given one forms with
$\sigma$; that attack reads messages and controls their delivery only
through the authenticated classical channel, so $\rho$ is an output state
that channel admits. Any measurement on the final state completes this
distinguisher, and the distinguishing advantage between two states is their
trace distance \cite[Sec.~III.B.1]{portmann2022security}, so
$\frac12\|\rho'-\rho\|_1\le\varepsilon_{\mathrm{auth}}$. By (ii),
$\frac12\|\rho-\mathcal R_{\mathrm{ideal}}(\rho)\|_1\le\epa+\eev$. The map
$\mathcal R_{\mathrm{ideal}}$ is a channel on the key registers
\cite[Sec.~3.3, Eqs.~(13), (17)]{tupkary2026decoy}, so it does not increase the
trace norm and
$\frac12\|\mathcal R_{\mathrm{ideal}}(\rho)-\mathcal R_{\mathrm{ideal}}(\rho')\|_1
\le\varepsilon_{\mathrm{auth}}$. The triangle inequality over $\rho'$, $\rho$,
$\mathcal R_{\mathrm{ideal}}(\rho)$ and $\mathcal R_{\mathrm{ideal}}(\rho')$
gives \eqref{eq:composed-auth}. The run consumes at most $k_{\mathrm{auth}}$
bits of pre-shared key (Assumption~\ref{ass:auth}\,(ii)) and, when both
parties output a key, produces an $\ell$-bit key, which is the accounting of
key consumed against key produced in \cite[Sec.~VII.D]{portmann2022security};
the net key gain is then $\ell-k_{\mathrm{auth}}$.
\end{proof}
